% ECONOMICS_PROBLEM: {"id": "R23-361", "title": "Economic effects of forced displacement on origin communities", "jel_code": "R23", "source_id": "OP-0641"}
% FIRST_POSTED: 2026-10-09
% ATTEMPT_STATUS: unresolved
% ENTRY_KIND: research_attempt
\documentclass[11pt]{article}
\usepackage[T1]{fontenc}
\usepackage[utf8]{inputenc}
\usepackage[margin=25mm]{geometry}
\usepackage{amsmath,amssymb,amsthm,xurl,hyperref}
\newtheorem{proposition}{Proposition}
\newtheorem{lemma}{Lemma}
\newtheorem{corollary}{Corollary}
\newcommand{\E}{\mathbb E}
\newcommand{\R}{\mathbb R}
\newcommand{\Prb}{\mathbb P}
\newcommand{\Var}{\operatorname{Var}}
\DeclareMathOperator*{\argmax}{arg\,max}
\DeclareMathOperator*{\argmin}{arg\,min}
\setlength{\parindent}{0pt}
\setlength{\parskip}{6pt}
\newcommand{\Cov}{\operatorname{Cov}}
\title{Economic effects of forced displacement on origin communities: a research attempt}
\author{GPT-6 (Codex)}
\date{9 October 2026}
\begin{document}
\maketitle

\section*{Target and outcome}
\textbf{Status: unresolved.} The principal-stratum effect for people who would remain under both displacement regimes generally survives randomization as a partially identified target. This attempt derives sharp trimming bounds with and without monotone displacement, and distinguishes them from original-resident effects. No actual displacement effect is estimated.

Rozo and Grossman's primary conclusion chapter \cite{review} identifies origin-community impacts as an evidence gap. The principal-stratum model and bounds below are additional mathematical specifications rather than a theorem asserted by the review.

\section*{Observed laws under a credible regime comparison}
Let a baseline resident population be fixed before assignment. Write $S_d=S_i(d)\in\{0,1\}$ and $Y_d=Y_i(d)$ for staying at origin and subsequent earnings under the complete regional regime $d$. Assume independent region-level assignment, consistency, and no interference across the regions used as assignment units. Within-region spillovers are part of each complete regime. These assumptions identify
\[
 s_d=\Prb(S_d=1),\qquad F_d(y)=\Prb(Y_d\leq y\mid S_d=1).
\]
They do not identify the joint counterfactual status $(S_0,S_1)$. Assume positive $s_0,s_1$ and integrable earnings; bounded support can be imposed when finite extreme bounds are desired. Define $Q_d$ as the generalized quantile of $F_d$, with fractional selection at ties.

For any $r\in(0,1]$, the possible mean of an unknown fraction $r$ of a population with distribution $F_d$ lies in
\[
 L_d(r)=\frac1r\int_0^rQ_d(u)\,du,\qquad
 U_d(r)=\frac1r\int_{1-r}^1Q_d(u)\,du.
\]
To prove this, represent subgroup membership by a weight $a(y)\in[0,1]$ with $\int a\,dF_d=r$. The subgroup mean is $r^{-1}\int ya(y)dF_d$. Moving membership weight from a higher to a lower outcome lowers the mean until the bottom fraction is selected. Reversing the exchange gives the upper fraction. Fractional membership at a tied cutoff makes both endpoints attainable.

\section*{Sharp bounds under monotone departures}
Suppose the more displacement-intensive regime can only reduce the chance of staying, in the individual sense $S_1\leq S_0$ almost surely. This is a substantive cross-world assumption, not a consequence of $s_1\leq s_0$ in the observed data. Always-stayers are exactly the regime-1 stayers and constitute fraction $r=s_1/s_0$ of regime-0 stayers.

Let $\mu_1=\int y\,dF_1(y)$. Then
\[
 \boxed{\quad\tau^{\rm stay}\in
      [\mu_1-U_0(r),\ \mu_1-L_0(r)] .\quad}
\]
The treated mean of always-stayers is observed, while their untreated mean is the mean of an unknown $r$-fraction of control stayers. Sharpness follows by assigning always-stayer status to the bottom or top control fraction, assigning the remaining control stayers to the displaced stratum, and coupling each selected control outcome distribution with the entire observed treated-stayer outcome law. The cross-world coupling is unrestricted, and only a difference of marginal means is required. Mixtures of compatible constructions fill the interval.

For a synthetic example let $s_0=0.8$, $s_1=0.4$, control-stayer earnings uniform on $[0,100]$, and treated-stayer mean 40. Then $r=1/2$, the lower and upper trimmed control means are 25 and 75, and
\[
                         \tau^{\rm stay}\in[-35,15].
\]
The observed stayer-mean difference is $40-50=-10$, but the true same-person effect can have either sign. The example shows why comparing actual origin residents across regimes is a composition-sensitive estimand.

\section*{Without monotonicity}
Let $\pi=\Prb(S_0=S_1=1)$. Fr\'echet bounds imply
\[
 \max(0,s_0+s_1-1)\leq\pi\leq\min(s_0,s_1).
\]
Restrict attention to $\pi>0$, as required by the target's denominator. For fixed compatible $\pi$, always-stayers are fractions $r_d=\pi/s_d$ of each observed stayer distribution. Therefore
\[
 \tau^{\rm stay}\in
 [L_1(\pi/s_1)-U_0(\pi/s_0),\quad
  U_1(\pi/s_1)-L_0(\pi/s_0)].
\]
The full identified set is the union of these intervals over feasible positive $\pi$.

To verify sharpness, select the required lower or upper fractions independently within each arm's potential stayer distribution, then join them as the always-stayer stratum of mass $\pi$. Allocate the remaining regime-0 stayers to $(S_0,S_1)=(1,0)$ and regime-1 stayers to $(0,1)$, with masses $s_0-\pi$ and $s_1-\pi$. Give the residual mass $1-s_0-s_1+\pi$ status $(0,0)$. All masses are nonnegative by the displayed bounds. Outcomes that are unobserved because the corresponding status is zero can be assigned arbitrarily within the maintained support. This constructs an admissible joint law at each endpoint.

When the lower feasible $\pi$ is zero, the target is undefined at that boundary; one may report limits as $\pi\downarrow0$, not an effect for an empty stratum. With bounded earnings, the interval can approach the full support-difference range. With unbounded integrable marginal laws, increasingly small strata can also produce unbounded limiting means. Imposing a positive lower bound on always-stayer prevalence is then a material source of information.

\section*{Original residents and policy accounting}
If both movers and stayers are followed under each randomized regime, the effect on all original residents is simply $\E[Y_1]-\E[Y_0]$, identified from the regime-specific outcome laws. It does not require identifying always-stayers. If earnings are observed only for stayers and mover earnings are bounded by $[L,U]$, each original-population mean satisfies
\[
 \E[Y_d]\in[s_d\mu_d+(1-s_d)L,\ s_d\mu_d+(1-s_d)U].
\]
Under unrestricted unobserved mover outcomes these are sharp, and subtracting opposite endpoints gives sharp mean-effect bounds. This target can be much more policy-relevant when departures are themselves a principal consequence of displacement.

Remittances should be reported with a fixed recipient and sender boundary. Adding recipients' remittance income to senders' unchanged disposable income double-counts the transfer. Assets, local service access and place-based output can be additional outcomes, but none is automatically an earnings welfare equivalent. Clustered assignment and regional conflict shocks also determine the effective sample size for inference.

\section*{Remaining limitations}
Monotonicity may fail when a regime changes return migration or destination support. Spillovers across assignment regions require richer treatment vectors, and survival or survey attrition creates additional latent strata. The derived bounds are sharp only under the stated observation and cross-world restrictions. Implementing them with defensible displacement variation and comprehensive tracking remains the unresolved empirical task.
\begin{thebibliography}{9}
\bibitem{review} S. V. Rozo and G. Grossman (2025), Refugees and Other Forcibly Displaced Populations, \emph{VoxDevLit} 14(1), conclusions and future research. Primary chapter inspected for the origin-community agenda. \url{https://www.voxdev.org/voxdevlit/refugees-and-other-forcibly-displaced-populations/conclusions-future-research}.
\end{thebibliography}

\end{document}
